\documentclass[11pt]{article}
\usepackage[T1]{fontenc}
\usepackage{lmodern,amsmath,amssymb,amsthm,mathtools,microtype}
\usepackage[margin=1in]{geometry}
\usepackage[colorlinks=true,linkcolor=blue,citecolor=blue,urlcolor=blue]{hyperref}
\usepackage{enumitem,booktabs}
\setlength{\emergencystretch}{1em}
\newtheorem{theorem}{Theorem}[section]
\newtheorem{lemma}[theorem]{Lemma}
\newtheorem{proposition}[theorem]{Proposition}
\newtheorem{corollary}[theorem]{Corollary}
\theoremstyle{remark}\newtheorem{remark}[theorem]{Remark}
\newcommand{\F}{\mathbb F}
\newcommand{\E}{\mathbb E}
\newcommand{\Prb}{\mathbb P}
\newcommand{\rank}{\operatorname{rank}}
\newcommand{\supp}{\operatorname{supp}}
\newcommand{\cc}{\operatorname{cc}}
\newcommand{\poly}{\operatorname{poly}}
\newcommand{\eps}{\varepsilon}
\newcommand{\calF}{\mathcal F}
\title{Research audit and proposed strengthenings\\\large Distributed vertex connectivity and vertex-deletion sketches}
\author{Working research note}
\date{October 28, 2026}
\begin{document}
\maketitle
\begin{abstract}
This note develops and checks several additions to the supplied draft \emph{Annihilator Sketches and the Complexity of Distributed Vertex Connectivity}: a constant-field public-randomness rank protocol with one-sided error; one-sided private-randomness decoding; polynomial-time construction of a genuinely sparse certificate; a randomized $\Omega(k)$ decision lower bound and an entropy lower bound for strong certificates; refined logarithmic factors in both upper bounds; an improved deterministic protocol near the complete-graph threshold; and an optimality observation for the generic linear deletion primitive. It also identifies errors and unfinished proof obligations in the draft. The arguments below underwent a separate self-audit and finite computations, but not an independent-agent review: no callable independent-agent capability was available. No claim of literature novelty is made for the derived additions.
\end{abstract}

\section{Conventions and summary}
All graphs are simple and undirected. A graph on more than $k$ vertices is $k$-vertex-connected if deletion of every set of \emph{fewer than} $k$ vertices leaves it connected. A strong certificate of order $k$ is a spanning subgraph $H\subseteq G$ satisfying
\[
 \cc(H-F)=\cc(G-F)\qquad\text{for every }F\subseteq V,\quad |F|<k,
\]
where equality means equality of the component partitions. Logarithms in asymptotic bounds may be taken in any fixed base.

For any fixed desired polynomial failure probability, the following bounds follow from the draft's primitive after the repairs described below:
\begin{align*}
 B_{\rm pub}&=O\!\left(\min\left\{n,\,k\log^2 n\log(en/k)\right\}\right),\\
 B_{\rm priv}&=O\!\left(\min\left\{n,\,\sqrt{nk\log(en/k)}\,\log n\right\}\right).
\end{align*}
The public protocol has a dynamic-streaming implementation using $O(nB_{\rm pub})$ bits. Both protocols can produce an $O(kn)$-edge strong certificate using polynomial-time postprocessing. Separately, the private messages support a one-sided-error decision decoder: it never accepts a graph that is not $k$-vertex-connected. The one-sided decoder described here tests all failure sets; polynomial-time one-sided decoding for variable $k$ is \emph{not} claimed.

An additional public-randomness \emph{one-sided decision} protocol uses
\[
 B_{\rm pub,1s}=O\!\left(\min\left\{n,\,
 \sqrt{n\log n}+k\log(en/k)\right\}\right)
\]
for a fixed polynomial error target. Unlike the certificate protocols, this new rank protocol has an all-failure-sets decoder, not a claimed polynomial-time global decoder. It supports simultaneous component-partition queries, but does not output a subgraph certificate. Section~\ref{sec:near-complete} further improves the deterministic decision bound when $k$ is very close to $n$.

For $1\le k\le n/6$, the decision lower bounds established below are
\[
 B_{\rm det}=\Omega(k+\sqrt n),\qquad
 B_{\rm rand}=\Omega(k+\log^3n),
\]
where the latter permits public randomness and constant error. The square-root term uses the supplied draft's Theorem 4 at $k=1$; the logarithmic term uses Yu's connectivity lower bound~\cite{Yu21}. For strong-certificate output, the $k$ term strengthens to $k\log(en/k)$ by an elementary information argument (Section~\ref{sec:certificate-lower}).

\section{One-sided private-randomness decoding}
The essential observation is that a rank calculation can certify connectivity even when sparse recovery could hallucinate edges.

\paragraph{Fixed edge coordinates and transformed sketches.}
Write $\mathcal E=\binom{[n]}2$ and $N=\binom n2$. Every incidence row and every aggregate incidence row is indexed by all $N$ original potential edges, including absent edges. AGM-style contraction sums rows; it does not replace this edge universe by the pairs of supernodes. Orient an edge $\{u,v\}$, $u<v$, with coefficient $+1$ at $u$ and $-1$ at $v$, and write $\epsilon_{v,u}=1$ for $v<u$ and $-1$ for $v>u$.

For clarity, the draft's raw neighbor sketch is
\[
 s(v)_q=\sum_{u\in N_G(v)}\epsilon_{v,u}a_u^q,
 \qquad 0\le q\le 2d+k.
\]
Its vertex-indexed Vandermonde map has $n$ columns. The \emph{transformed} edge-indexed map below is a different matrix with $N$ columns. Write
\[
 P_F(X)=\prod_{z\in F}(X-a_z)
       =\sum_{\ell=0}^{|F|}\beta_\ell X^\ell,
 \qquad t=2d+k-|F|+1.
\]
For $v\notin F$, the coordinator computes the $t$-coordinate row
\begin{align*}
 \widetilde s_F(v)_j
 &=a_v^jP_F(a_v)\sum_{\ell=0}^{|F|}\beta_\ell s(v)_{j+\ell}\\
 &=\sum_{u\in N_G(v)\setminus F}
   \epsilon_{v,u}P_F(a_v)P_F(a_u)(a_va_u)^j,
 \qquad 0\le j<t.
\end{align*}
This fixes the sign convention consistently with the draft's displayed raw sketch. The transformation uses only $F$, the public labels, and the stored sketch, without knowledge of the graph. The value of $t$ above retains all coordinates in the draft's construction; the rank argument below also applies to a shorter retained Vandermonde map whenever it has the stated column-rank guarantee.

\begin{lemma}[Deterministic rank certificate]\label{lem:rank}
Fix a failure set $F$ with $|F|\le k$, and let $J\subseteq G-F$ be any known spanning subgraph with components $C_1,\ldots,C_r$. Aggregate the transformed deterministic sketches over each $C_i$, and let $M\in\F_Q^{r\times t}$ have rows $M_{i,\cdot}=\sum_{v\in C_i}\widetilde s_F(v)$. Write $c=|\cc(G-F)|$. Then
\[
 \rank M\le r-c.
\]
If at most $d$ edges of $G-F$ join different components of $J$, and the transformed sketch has the draft's $d$-sparse-recovery Vandermonde guarantee, then equality holds.
\end{lemma}
\begin{proof}
Let $\mathcal B_F\in\F_Q^{(n-|F|)\times N}$ be the signed incidence matrix with the fixed column set $\mathcal E$: a surviving edge of $G-F$ has the two oriented nonzero entries just specified, while an absent edge or an edge incident to $F$ has a zero column. Let $R\in\F_Q^{r\times(n-|F|)}$ have $R_{i,v}=1$ if $v\in C_i$ and zero otherwise. Thus $B=R\mathcal B_F\in\F_Q^{r\times N}$ is the aggregate incidence matrix. Edges internal to a component cancel, leaving their original columns equal to zero; parallel edges between supernodes keep their distinct original columns.

For every potential edge $e=\{u,v\}\in\mathcal E$, set
\[
 W_{e,e}=P_F(a_u)P_F(a_v),\qquad
 \Psi_{j,e}=(a_ua_v)^j\quad(0\le j<t),
\]
with all off-diagonal entries of $W$ zero. In particular,
\[
 W\in\F_Q^{N\times N},\qquad
 \Psi\in\F_Q^{t\times N},\qquad
 \underbrace{M}_{r\times t}
 =\underbrace{B}_{r\times N}
  \underbrace{W}_{N\times N}
  \underbrace{\Psi^{\mathsf T}}_{N\times t}.
\]
The diagonal entry of $W$ is zero when an endpoint fails and is nonzero when both endpoints survive. Thus $W$ need not be invertible on the full edge universe. Distinct original edges retain distinct product labels $a_ua_v$, under the draft's field-label assumption. Viewing the nonzero columns of $B$ as a contracted multigraph is useful for its rank: it has $c$ components, so $\rank B=r-c$. The full-size factorization therefore gives $\rank M\le r-c$ for every graph and every sampling outcome.

For the equality proof only, let $S\subseteq\mathcal E$ consist of the surviving edges of $G-F$ joining different components of $J$, and put $m=|S|$. Define the analytic restrictions
\[
 B_S=B_{[:,S]}\in\F_Q^{r\times m},\quad
 W_S=W_{[S,S]}\in\F_Q^{m\times m},\quad
 \Psi_S=\Psi_{[:,S]}\in\F_Q^{t\times m}.
\]
All columns of $B$ outside $S$ are zero, so $\rank B_S=\rank B$ and $M=B_SW_S\Psi_S^{\mathsf T}$. These smaller matrices are proof-only restrictions, not new protocol coordinates and not matrices obtained by replacing original edges with supernode pairs. If $m=0$, then $B=M=0$ and equality is immediate. If $1\le m\le d$, the Vandermonde guarantee implies that $\Psi_S$ has full column rank; only this restricted matrix, not necessarily $\Psi$, has that property. Also $W_S$ is invertible. Hence $W_S\Psi_S^{\mathsf T}\in\F_Q^{m\times t}$ has a right inverse $T_S\in\F_Q^{t\times m}$. Multiplying gives $MT_S=B_S$, so $\rank M\ge\rank B_S=r-c$. Together with the upper bound, this proves equality.
\end{proof}

\begin{corollary}[One-sided connectivity test]
Output ``connected'' exactly when $\rank M=r-1$. This answer is always sound. It is complete whenever the inter-component edge bound in Lemma~\ref{lem:rank} holds.
\end{corollary}

In the private protocol, take $J=H-F$, where $H$ is the union of the actual sampled edges sent by the vertices. For all $|F|<k$, apply the test and accept $k$-vertex connectivity only if every test accepts. A non-$k$-connected graph is rejected for every random outcome, because one of its genuine separators supplies a rank deficiency. On the simultaneous sparse-boundary event used in the draft, all tests are exact. Thus the communication bound is unchanged, but the decision error becomes one-sided.

There is a stronger partition interpretation. Let $K=\ker M^{\mathsf T}\subseteq\F_Q^r$, and declare
\[
 i\sim j\quad\Longleftrightarrow\quad x_i=x_j\text{ for every }x\in K.
\]
The indicator of each genuine component of $G-F$ belongs to $K$, so this equivalence relation never merges distinct genuine components. When the sparse-boundary event holds, the kernel has dimension $c$ and is exactly the span of these component indicators. Consequently the computed partition is then exact. This gives one-sided component recovery, but not a guarantee that an arbitrary output forest consists of genuine edges.

For $k=1$, this is a polynomial-time rank calculation with the usual private-randomness communication. HKTZZ discusses one-sided-error connectivity near the end of its Section 3~\cite{HKTZZ19}, but that sentence also uses the term Las Vegas. The result here is a one-sided Monte Carlo protocol, not a zero-error Las Vegas protocol. It should not be advertised as resolving an intended zero-error question. Whether the rank observation already appears in subsequent work has not been established by the present literature check.

\section{A stronger public-randomness decision bound over a constant-size field}\label{sec:constant-rank}
The alternate fingerprint construction in the draft can do more than verify a candidate sparse vector. With enough independent coordinates, it preserves the rank of the entire contracted incidence space. Crucially, this application works over the \emph{constant-size field} $\F_5$; the large field needed for Vandermonde sparse recovery is unnecessary here.

\begin{theorem}[One-sided public-randomness decision]\label{thm:constant-rank}
For $1\le k<n$ and $0<\delta\le1/2$, there is a simultaneous-message public-randomness protocol for $k$-vertex connectivity with maximum message length
\[
 O\!\left(\min\left\{n,\,
 \sqrt{n\log(2n)}+k\log(en/k)+\log(1/\delta)\right\}\right)
\]
bits. It never accepts a graph that is not $k$-vertex-connected, and accepts every $k$-vertex-connected graph with probability at least $1-\delta$.

More strongly, its messages support component-partition queries for every $|F|<k$: the answer always refines the true partition of $G-F$, and with probability at least $1-\delta$ all the answers are exact simultaneously. A query takes polynomial time. The global decision decoder enumerates failure sets, and is \emph{not claimed to be polynomial-time for variable $k$}. The protocol does not output a subgraph certificate.
\end{theorem}

\subsection{A deterministic preliminary subgraph}
Fix an integer $b\ge1$, and put $s=b+k-1$. Each vertex sends a deterministically chosen set of $\min\{s,\deg(v)\}$ neighbors, for example the lowest-labeled such neighbors. Let $J$ be the union of these genuine edges. This choice is independent of the subsequent sketching randomness.

For a fixed $|F|<k$, let $C_1,\ldots,C_r$ be the components of $J-F$, and let $c$ be the number of components of $G-F$. Call $C_i$ active if it is not an entire component of $G-F$. An active $C_i$ has an edge $vu$ of $G-F$ leaving it. Since this edge was not contributed by $v$, the vertex $v$ contributed $s$ neighbors; at least $s-(k-1)=b$ survive and belong to $C_i$. Thus each active component has at least $b+1$ vertices. In particular,
\begin{equation}\label{eq:rank-budget}
 h:=r-c\le \#\{\text{active components}\}\le H:=\left\lceil\frac{n}{b+1}\right\rceil.
\end{equation}
This holds for every failure set and every graph, not just for $k$-connected graphs.

\subsection{The constant-field messages and their transformation}
Set
\[
 M_k=\sum_{j=0}^{k-1}\binom nj,\qquad
 \lambda=\left\lceil\log_2(M_k/\delta)\right\rceil+1,\qquad
 D=3H+\lambda,\qquad m=D+k-1.
\]
Using public randomness, choose a matrix $\Phi\in\F_5^{m\times n}$ with independent uniform entries from the \emph{whole} field. In addition to its selected neighbors, vertex $v$ sends
\[
 f(v)=\sum_{u\in N(v)}\epsilon_{v,u}\Phi_u\in\F_5^m,
 \qquad
 \epsilon_{v,u}=\begin{cases}1,&v<u,\\-1,&v>u,\end{cases}
\]
where $\Phi_u$ is column $u$ of $\Phi$. These are $m$ field elements and hence $O(m)$ bits, not $O(m\log n)$ bits.

For a query $F$, deterministically choose $D$ linearly independent vectors
\[
 p_1,\ldots,p_D\in\ker(\Phi_F^{\mathsf T}).
\]
Here $\Phi_F\in\F_5^{m\times|F|}$ and each $p_t\in\F_5^m$. There is room for them because the kernel dimension is at least $m-|F|\ge D$. Set $z_t(v)=p_t^{\mathsf T}\Phi_v$. For each component $C_i$ of $J-F$, form the matrix $M\in\F_5^{r\times D}$ with entries
\[
 M_{i,t}=\sum_{v\in C_i}z_t(v)\,p_t^{\mathsf T}f(v),\qquad t\in[D].
\]
The draft's cancellation identity, with a consistent orientation, gives
\begin{equation}\label{eq:constant-transform}
 M_{i,t}=\sum_{\substack{v\in C_i,\ u\notin C_i\cup F\\uv\in E(G)}}
 \epsilon_{v,u}z_t(v)z_t(u).
\end{equation}
The same fixed-edge convention applies here. Let $B\in\F_5^{r\times N}$ be the aggregate incidence matrix defined as in Section~2, and let $K\in\F_5^{N\times D}$ have $K_{\{u,v\},t}=z_t(u)z_t(v)$. Equation~\eqref{eq:constant-transform} says $M=BK$, a product of dimensions $(r\times N)(N\times D)=r\times D$. The vertex-indexed random matrix $\Phi\in\F_5^{m\times n}$ is not an edge-incidence matrix and does not change the ambient edge universe. Thus every column of $M$ lies in $S=\operatorname{im}B\subseteq\F_5^r$; viewing the nonzero columns as a contracted multigraph shows that $\dim S=r-c$. In particular,
\[
 \rank M\le\dim S=r-c=h
\]
for \emph{every} outcome of the randomness. Testing $\rank M=r-1$ can therefore never falsely certify connectivity.

\subsection{Why the rank is full with high probability}
Condition on any value of $\Phi_F$. The selected vectors $p_1,\ldots,p_D$ are now fixed, and their linear independence implies that
\[
 (z_1(v),\ldots,z_D(v))
\]
is uniform in $\F_5^D$ for each $v\notin F$. These vectors are independent across surviving vertices, and all their coordinates are independent.

Take a nonzero linear functional on $S$, represented by $x\in\F_5^r$ that does not annihilate $S$. Write $\pi(v)$ for the index of the component containing $v$. On column $t$ of $M$, this functional is
\begin{equation}\label{eq:cut-poly}
 x^{\mathsf T}M_{\cdot,t}=
 \sum_{\substack{\{u,v\}\in E(G-F)\\u<v}}
 \bigl(x_{\pi(u)}-x_{\pi(v)}\bigr)z_t(u)z_t(v).
\end{equation}
This is a nonzero polynomial of total degree two: some contracted edge has endpoints with different $x$ values, and distinct original edges give distinct monomials. The fact that $G$ is simple is used here. Schwartz--Zippel gives
\[
 \Prb[x^{\mathsf T}M_{\cdot,t}=0\mid\Phi_F]\le 2/5.
\]
Different coordinates $t$ use independent random variables. There are fewer than $5^h$ nonzero linear functionals on $S$. Consequently,
\begin{align*}
 \Prb[\rank M<h\mid\Phi_F]
 &\le 5^h(2/5)^D\\
 &\le (8/25)^H(2/5)^\lambda
 \le 2^{-\lambda}.
\end{align*}
The middle inequality uses $h\le H$ and $D=3H+\lambda$. For $h=0$, the rank is always exact and no union bound is needed. Averaging over $\Phi_F$ and union-bounding over the $M_k$ failure sets proves simultaneous exactness with probability at least $1-\delta$.

The kernel-equivalence rule from Section~2 recovers the component partition when the rank is exact, and never merges distinct true components otherwise. This proves the correctness and error-direction claims in Theorem~\ref{thm:constant-rank}.

\subsection{Bit complexity and the importance of set encoding}
A subset of at most $s$ labels from $[n]$ can be encoded in
\[
 O\!\left(s\log(en/s)\right)
\]
bits when $s\le n$, with full adjacency transmission available as a fallback. An enumerative encoding suffices and has polynomial-time encoding and decoding. Using a list of labels would be valid but would lose part of the improvement.

The raw field-vector cost is
\[
 O(m)=O\!\left(k+\frac nb+\log M_k+\log(1/\delta)\right).
\]
Since $\log M_k=O(k\log(en/k))$, the total maximum message length is
\[
 O\!\left(
 b\log(en/b)+\frac nb+k\log(en/k)+\log(1/\delta)
 \right).
\]
Taking $b=\lceil\sqrt{n/\log(2n)}\rceil$ proves the theorem, with the full-graph protocol providing the $n$-bit cap.

\paragraph{What is and is not improved.}
At $k=\Theta(\sqrt n)$ and polynomially small error, this is $O(\sqrt n\log n)$ bits per vertex, versus the draft's $O(\sqrt n\log^3n)$ public-randomness bound. It also has one-sided error. For small $k$, the ordinary public-randomness forest protocol is better in communication, and unlike this rank protocol it gives an explicit sparse subgraph. The new result is therefore an additional tradeoff, not a replacement for every headline theorem.

\paragraph{A weaker but explicit streaming consequence.}
The same constant-field rank construction is linear in the graph and its entire matrix $\Phi$ costs $O(nm)$ bits to store. The preliminary subgraph can be constructed by the certified sparse-recovery initialization in Section~\ref{sec:public}, replacing $2k$ by $s=b+k-1$. Using capacity $O(s+\log n)$ at geometric levels costs $O((s+\log n)\log^2n)$ bits per vertex. Exact sampled-degree counts certify that every recovered edge is genuine; abort if some high-degree vertex does not supply $s$ neighbors. This abort event has polynomially small probability. The initialization and $\Phi$ use independent randomness.

Taking $b=\lceil\sqrt n/\log(2n)\rceil$ yields, for any fixed polynomial error target, a one-pass dynamic-streaming decision algorithm with one-sided error and space
\[
 O\!\left(\min\left\{n^2,\,
 n\bigl(\sqrt n\log n+k\log^2n\bigr)\right\}\right)
\]
bits. The same exponential-in-$k$ postprocessing caveat applies. This streaming bound is distinct from, and usually weaker than, the distributed message bound multiplied by $n$, because the streamed initialization cannot simply choose and retain arbitrary neighbors for free.

\section{Polynomial-time construction of a sparse strong certificate}
The draft enumerates failure sets and unions spanning forests. Even correct forests do not by themselves bound the size of that union. The following black-box reduction replaces that enumeration.

Assume access to an oracle $\mathcal O(F)$ that returns a spanning forest of $G-F$ for any $|F|<k$. In the application, the oracle is the coordinator's sketch decoder. Work on the single event that all such oracle answers are correct.

\begin{theorem}\label{thm:sparsifier}
Put $C_k=kn-\binom{k+1}{2}$. Using at most $\binom n2+C_k$ oracle calls and polynomial additional computation, one can construct a strong certificate $H\subseteq G$ with at most $C_k$ edges. The additional graph storage is $O(kn\log n)$ bits, or $O(n^2)$ bits using an adjacency representation.
\end{theorem}
\begin{proof}
Start with $H=(V,\varnothing)$ and process unordered pairs $\{u,v\}$ in a fixed order. For the current pair, repeat the following procedure.

If $uv\in E(H)$, finish this pair. Otherwise compute the minimum size of a $u$--$v$ vertex separator in $H$, excluding the endpoints. If this size is at least $k$, finish the pair. Otherwise choose such a separator $F$ with $|F|<k$ and query $\mathcal O(F)$.

If the returned forest separates $u$ and $v$, finish the pair. This establishes $uv\notin E(G)$, since an actual edge $uv$ would survive deletion of $F$. If the forest connects them, its $u$--$v$ path contains an edge $xy$ crossing components of $H-F$. Add one such edge to $H$, and repeat on the same pair. This is a genuine edge of $G$, and before insertion its endpoints were separated by fewer than $k$ vertices.

We next bound the number of insertions. The insertion order is a $k$-construction order: immediately before each edge is inserted, its endpoints have at most $k-1$ internally vertex-disjoint paths. Theorem 6 of Sorkin--Steger--Zenklusen~\cite{SSZ09} therefore bounds the number of inserted edges by $C_k=kn-\binom{k+1}{2}$.

A coarser $2kn$ bound follows using only the Mader-based argument reviewed in Appendix A of Assadi--Shah~\cite{AS23}. The resulting graph cannot contain a $(k+1)$-vertex-connected subgraph: immediately before the last edge of such a subgraph was inserted, its endpoints would already have $k$ internally disjoint paths. Mader's density theorem then gives $2kn$. The sharper constant in the present theorem is an application of the existing extremal result, not a new extremal theorem.

A query either adds an edge or finishes the current pair. Thus the number of calls is at most $C_k+\binom n2$.

At termination, every edge $uv\in E(G)\setminus E(H)$ has at least $k$ internally vertex-disjoint paths in $H$. It cannot have been dismissed by a negative oracle answer, and the other dismissal condition is monotone as $H$ grows. For any $|F|<k$, every surviving omitted edge therefore has endpoints connected in $H-F$. Replacing omitted edges along paths of $G-F$ proves that $H-F$ and $G-F$ have the same component partition.
\end{proof}

\paragraph{Adaptivity.}
The failure sets depend on previous oracle answers. A guarantee for each individually fixed $F$ is not sufficient by itself. The draft's intended simultaneous correctness event over \emph{all} $F$ makes the argument valid without fresh sketches. This reduction retains the original two-sided-error guarantee; it is not automatically a polynomial-time implementation of the one-sided rank decoder.

\paragraph{Decision only.}
A simpler variant repeatedly finds any separator $F$ of the current $H$ with $|F|<k$. If $G-F$ is disconnected, reject; otherwise add a genuine forest edge crossing components of $H-F$. Accept once $H$ becomes $k$-connected. The same insertion argument bounds this by $C_k+1$ forest queries.

\section{A linear-in-k randomized decision lower bound}
A streaming lower bound does not automatically imply the analogous node-sketching lower bound. Here is a direct reduction that does.

\begin{lemma}[Two cliques and a bipartite graph]\label{lem:cliques}
Let $L,R$ each induce a clique of size at least $k+1$, and let $D$ be the bipartite graph of edges between them. Then the resulting graph is $k$-vertex-connected if and only if $D$ has a matching of size at least $k$.
\end{lemma}
\begin{proof}
A matching of size $k$ leaves a cross-edge after deletion of fewer than $k$ vertices; both clique remainders are nonempty and connected. Conversely, if the maximum matching has size less than $k$, the bipartite matching--cover theorem supplies a vertex cover of $D$ of size less than $k$. Deleting it removes all cross-edges while leaving a vertex in each clique.
\end{proof}

\begin{theorem}\label{thm:index}
For $3\le k$ and $n\ge2k+2$, any randomized simultaneous-message protocol deciding $k$-vertex connectivity with error at most $1/3$ has worst-case maximum message length $\Omega(k)$, even with public randomness.
\end{theorem}
\begin{proof}
First suppose $k=2s+1$. Alice has an arbitrary matrix $X\in\{0,1\}^{s\times s}$; Bob has an index $(i,j)$ and wants $X_{ij}$.

Use two cliques
\[
 L=A\cup B'\cup\{p,\ell\},\qquad
 R=B\cup A'\cup\{q,r\},
\]
where $A=\{a_1,\ldots,a_s\}$ and $B,A',B'$ have corresponding labels and the same size. Each clique has $2s+2=k+1$ vertices. The cross-edges are
\[
 a_ta'_t,\quad b'_tb_t\quad(t\in[s]);\qquad
 a_tb_u\quad(X_{tu}=1);\qquad
 pa'_i,\quad b'_jq.
\]
The buffer vertices $\ell,r$ have no cross-edges.

Any cross-matching of size $2s+1$ must match every nonbuffer vertex. Thus $p$ must match $a'_i$ and $q$ must match $b'_j$. Every other $a'_t$ forces $a_ta'_t$, and every other $b'_u$ forces $b'_ub_u$. The only remaining unmatched vertices are $a_i,b_j$, which can be matched exactly when $X_{ij}=1$. Conversely, these forced edges plus $a_ib_j$ give the desired matching when that bit is one. Lemma~\ref{lem:cliques} proves the reduction.

Only vertices in $A\cup B$ have neighborhoods depending on $X$, and none of their neighborhoods depends on $(i,j)$. Alice sends their $2s$ sketch messages. Bob can simulate all remaining vertices, whose neighborhoods depend only on his index and the fixed gadget. A protocol with maximum message length $B_0$ would therefore give a one-way indexing protocol using $O(s(B_0+1))$ bits.

For completeness, let $X$ be uniform and let $R_0$ contain all random coins. If $M$ is Alice's message tuple, the decoder for every bit has error at most $1/3$. Therefore
\[
 H(X\mid M,R_0)\le \sum_{t,u} H(X_{tu}\mid M,R_0)
 \le s^2 h_2(1/3),
\]
so $I(X;M\mid R_0)\ge(1-h_2(1/3))s^2$. Since each of the $2s$ messages has fewer than $2^{B_0+1}$ possibilities, $H(M\mid R_0)\le2s(B_0+1)$. Hence $B_0=\Omega(s)=\Omega(k)$ asymptotically; bounded $k$ is covered by a constant lower bound.

For even $k$, start with the gadget for $k-1$ and add one forced cross-pair, one endpoint in each clique, with no other cross-edges incident to the new endpoints. It raises the maximum cross-matching size by exactly one. For larger $n$, add padding vertices to either clique with no cross-edges. The same characterization remains valid.
\end{proof}

\begin{lemma}[Universal-vertex lifting]\label{lem:join}
Let $G$ have $m>j$ vertices and let $G^+=G\vee K_t$, the graph obtained by adjoining $t$ universal vertices. Then $G^+$ is $(j+t)$-vertex-connected exactly when $G$ is $j$-vertex-connected. A distributed protocol on $G^+$ induces one on $G$ with the same bound on maximum message length.
\end{lemma}
\begin{proof}
If a universal vertex survives deletion, the remainder is connected. Otherwise all $t$ universal vertices have been deleted, leaving fewer than $j$ deletions in $G$. The reverse implication follows by adjoining all universal vertices to a separator in $G$. The new vertices have fixed, known neighborhoods, so the coordinator can simulate their messages; each original vertex adds the same fixed set of new neighbors to its input.
\end{proof}

Apply Lemma~\ref{lem:join} with $j=1$ and $t=k-1$. The draft's $\Omega(\sqrt m)$ deterministic connectivity lower bound gives
\[
 B_{\rm det}=\Omega\bigl(\sqrt{n-k+1}\bigr).
\]
Combining with Theorem~\ref{thm:index} establishes the abstract's $\Omega(k+\sqrt n)$ for $k\le n/6$, indeed for $k\le n/2-1$ up to harmless endpoint cases. Applying the same lifting to Yu's randomized lower bound~\cite{Yu21} gives $\Omega(k+\log^3 n)$ in this range (constant-error amplification adjusts $1/3$ to $1/4$).

More generally, using a base parameter $j\le\min\{k,n-k-2\}$ in Theorem~\ref{thm:index} and adding $k-j$ universal vertices yields, with bounded exceptional cases absorbed into constants,
\[
 B_{\rm det}=\Omega\left(\min\{k,n-k\}+\sqrt{n-k+1}\right),
\]
and the analogous randomized bound with $\log^3(n-k+1)$ in place of the square root. These are bounds for \emph{exact threshold decision}, not for the promise problem of distinguishing disconnected from $k$-connected graphs.

\section{A stronger information lower bound for strong certificates}\label{sec:certificate-lower}
There is a further lower bound that applies to the draft's \emph{certificate output}, not merely to threshold decision. This distinction matters: a decision protocol need not identify any actual edges.

\begin{theorem}\label{thm:certificate-lower}
For $1\le k<n$, every randomized simultaneous-message protocol that, with probability at least $2/3$, outputs a strong certificate of order $k$ has maximum message length
\[
 \Omega\!\left(k\log(en/k)\right)
\]
bits per vertex, even with public randomness. Combining this with the ordinary connectivity lower bounds gives
\[
 B_{\rm rand,strong}=\Omega\!\left(k\log(en/k)+\log^3 n\right),\qquad
 B_{\rm det,strong}=\Omega\!\left(k\log(en/k)+\sqrt n\right).
\]
The square-root term uses the supplied draft's deterministic connectivity theorem.
\end{theorem}
\begin{proof}
First observe that every edge of a graph of maximum degree at most $k$ is indispensable in a strong certificate. Indeed, suppose $uv\in E(G)$ is omitted from $H\subseteq G$. Delete
\[
 F=N_G(v)\setminus\{u\}.
\]
This uses at most $k-1$ vertices and does not delete $u$ or $v$. The vertex $v$ is isolated in $H-F$ but is joined to $u$ in $G-F$. Hence $H$ is not a strong certificate. Consequently, on every graph of maximum degree at most $k$, the only possible strong certificate is $G$ itself.

There are $2^{\Omega(nk\log(en/k))}$ labeled graphs of maximum degree at most $k$. Here is an elementary count. Let $N=\lfloor n/2\rfloor$ and use two labeled parts of size $N$. For $k\le N/e^2$, take the union of an ordered $k$-tuple of perfect matchings. There are $(N!)^k$ tuples. A resulting graph has degree at most $k$, and is the image of at most $k^{kN}$ tuples: for each left vertex and each of the $k$ matching positions, there are at most $k$ choices for its neighbor. Thus the number of distinct graphs is at least
\[
 \frac{(N!)^k}{k^{kN}}
 \ge \left(\frac{N}{ek}\right)^{kN}
 =2^{\Omega(nk\log(en/k))}.
\]
For $k>N/e^2$, set $q=\min\{k,N\}$ and take arbitrary subgraphs of a fixed $K_{q,q}$, with the other vertices isolated. There are $2^{q^2}=2^{\Omega(n^2)}$ such graphs, which is more than sufficient because $k\log(en/k)=O(n)$. Bounded exceptional values of $n$ are absorbed into constants.

Choose $G$ uniformly from this family and reveal all random coins $R$ as public randomness. The coordinator must reconstruct $G$ with error at most $1/3$. By the elementary decoding-entropy bound,
\[
 H(G\mid\text{transcript},R)
 \le 1+\tfrac13\log_2|\mathcal G|.
\]
Thus the transcript carries $\Omega(nk\log(en/k))$ bits of information. If each message has length at most $B$, its alphabet has size less than $2^{B+1}$, so the transcript has entropy at most $n(B+1)$. This proves the claimed lower bound.

Finally, any strong certificate preserves the component partition with no deletions, so computing one solves ordinary connectivity. Yu's randomized connectivity lower bound~\cite{Yu21} and the supplied draft's deterministic lower bound give the remaining terms. Combining two lower bounds by taking their maximum is equivalent, up to a factor of two, to adding them.
\end{proof}

\paragraph{An extreme decision--certificate separation.}
At $k=n-1$, exact threshold decision requires only one bit per vertex (Section~\ref{sec:near-complete}), whereas strong-certificate construction requires $\Theta(n)$ bits per vertex. In fact, a strong certificate must preserve every two-vertex induced subgraph, and hence must equal $G$. The simple decision protocol should therefore not be mistaken for a certificate protocol. Conversely, this certificate lower bound does not apply to the rank protocol in Section~\ref{sec:constant-rank}, whose output is a partition oracle rather than a subgraph.

\section{A deterministic improvement near the complete-graph threshold}\label{sec:near-complete}
The draft falls back to $n$ adjacency bits per vertex for large $k$. For \emph{decision}, a complement encoding can be much smaller.

\begin{proposition}
Let $t=n-1-k$. There is a deterministic decision protocol with maximum message length at most
\[
 1+\left\lceil\log_2\left(\sum_{j=0}^{t}\binom{n-1}{j}\right)\right\rceil
\]
bits. For $1\le t\le(n-1)/2$, this is $O(1+t\log(en/t))$. For $k=n-1$ the problem needs only one bit per vertex. For $k=n-2\ge1$, one bit per vertex also suffices.
\end{proposition}
\begin{proof}
A vertex of degree less than $k$ sends a rejection flag, which is sound because $k$-vertex connectivity implies minimum degree at least $k$. Every other vertex has at most $t$ nonneighbors and sends an encoding of that nonneighbor set. If no vertex rejects, the coordinator reconstructs the graph exactly and decides the problem. The number of possible sets is the sum in the statement.

For $k=n-1$, the graph must be complete, so each vertex only needs to indicate whether its degree is $n-1$. For $k=n-2$, minimum degree at least $n-2$ is both necessary and sufficient. Indeed, the complement then has maximum degree at most one. After deleting fewer than $n-2$ vertices, at least three vertices survive, and a complete graph on at least three vertices minus a matching is connected. Thus each vertex sends just the bit indicating whether its degree is at least $n-2$.
\end{proof}
This observation concerns exact threshold decision. On a rejected graph it does not supply a strong certificate of the original graph. It also explains why a lower bound of $\Omega(k)$ cannot hold uniformly all the way to $k=n-1$.

\section{Sharper private-randomness parameters}\label{sec:private}
Let
\[
 M_k=\sum_{j=0}^{k-1}\binom nj,\qquad
 \lambda=\left\lceil\log_2(M_k/\delta)\right\rceil+1.
\]
For $\delta\le1/2$, $\lambda\ge k$. A sampling budget and recovery capacity of order $\sqrt{n\lambda}$ suffice, with full-neighborhood transmission used when that is cheaper in bits.

Here are the necessary details of the tail argument. Let each vertex select $s$ distinct uniform neighbors, or all its neighbors if its degree is smaller. For a fixed $F$, restrict a uniform ordering of each vertex's original neighborhood to its surviving neighbors. This gives an independent uniform ordering in $G-F$. The surviving sample contains a uniform $\sigma$-out sample with $\sigma=\lfloor s/2\rfloor$, provided $s\ge2k$. This is a coupling statement, not merely a lower bound on the number of surviving sampled neighbors.

HKTZZ proves that a uniform $t$-out sample has expected $O(n/t)$ inter-component edges for $t\ge c_0\log n$~\cite{HKTZZ19}. Couple a $\sigma$-out sample to the union of $\lambda$ independent $\lfloor\sigma/\lambda\rfloor$-out samples, padding the union uniformly if necessary. Each trial has probability at most $1/2$ of leaving more than $C n\lambda/\sigma$ inter-component edges. The union has no more inter-component edges than any trial. Consequently
\[
 \Prb\left[\text{inter-component edges}>C n\lambda/\sigma\right]\le2^{-\lambda},
\]
provided $\sigma/\lambda\ge C_0\log n$. This side condition is important: it is implicit in the splitting proof printed for HKTZZ's Corollary 2.22.

Take $s=A\sqrt{n\lambda}$ and recovery capacity $d=A'\sqrt{n\lambda}$, for suitable constants. In the nontrivial regime $\lambda=O(n/\log^2n)$, the side condition holds after adjusting constants. Outside this regime, transmitting $n$ adjacency bits per vertex already fits the claimed bound. A union bound over $M_k$ failure sets now gives failure at most $\delta$. Thus
\[
 B_{\rm priv}=O\left(\min\{n,\sqrt{n\lambda}\log n\}\right).
\]
For $\delta=n^{-c}$ with fixed $c$, the standard binomial estimate gives $\lambda=O(k\log(en/k))$ (and $\lambda=O(n)$ near $k=n$). This proves the private bound in Section 1. The improvement is a logarithmic refinement, most useful when $k$ is close to $n$ up to polylogarithmic factors; it is not a better polynomial dependence on $k$ or $n$.

\section{Fewer public-randomness merging rounds}\label{sec:public}
The public protocol can first obtain a genuine subgraph $J\subseteq G$ such that every vertex contributes either all its neighbors or at least $2k$ distinct neighbors. In the distributed model, sending any $\min\{2k,\deg(v)\}$ incident edges costs $O(k\log n)$ bits per vertex and gives this property deterministically. Randomness is not needed for this initialization.

\begin{lemma}[Large unfinished components]
For every $|F|<k$, each component of $J-F$ that is not a whole component of $G-F$ contains at least $k+2$ vertices.
\end{lemma}
\begin{proof}
There is an edge $vu$ of $G-F$ leaving that component. It is absent from $J$, so $v$ did not contribute all its neighbors and contributed at least $2k$ distinct neighbors. At most $k-1$ were deleted. All remaining contributed neighbors lie in the same component of $J-F$ as $v$; there are at least $k+1$ of them.
\end{proof}

Thus the initial excess number of blobs over true connected components is at most $n/(k+2)$, simultaneously for all failure sets. We need only
\[
 R=\Theta(\log(en/k))
\]
merging rounds, rather than $\Theta(\log n)$, provided we prove the appropriate high-probability progress bound.

\subsection{A rigorous progress lemma without independence between blobs}
Use recovery capacity $d=A k$, with a sufficiently large constant $A$ depending on the desired final polynomial error exponent. For an active blob, the sampling argument in the draft gives failure at most
\[
 p\le2e^{-d/12}.
\]
This follows by choosing a level whose expected sampled support lies between $d/4$ and $d/2$; blobs with at most $d$ boundary edges are handled at level zero.

Condition on the previous rounds of an ideal execution that discards incorrectly recovered vectors. Fresh randomness for the next round is independent of that history. If $a$ blobs are active, let $Z$ be the number failing to find an outgoing edge. Then $\E Z\le pa$, regardless of dependence between blobs. Markov's inequality gives $\Prb[Z>a/2]\le2p$. Otherwise at least $a/2$ blobs select outgoing edges. Merging these selections reduces the blob count by at least $a/4$.

If $X$ is the excess number of blobs over true components, then $a\ge X$. Therefore a good round sends $X$ to at most $3X/4$. The conditional probability of a bad round is at most $\beta=4e^{-d/12}$. For any specified set of $t$ rounds, the probability that they are all bad is at most $\beta^t$, by iterated conditioning. Hence
\[
 \Prb[\text{at least }R/2\text{ bad rounds}]
 \le2^R\beta^{R/2}=\exp(-\Omega(kR)).
\]
Choose the constant in $R$ so that $R/2$ good rounds eliminate an initial excess of at most $n/(k+2)$. Since
\[
 \log M_k=O(k\log(en/k)),
\]
the last probability is small enough for the union bound over all failure sets, and also for any fixed desired polynomial error bound. Fingerprint collisions are bounded separately along the ideal executions, so this proof never conditions on a future ``no hallucinations'' event.

For completeness, choose each fingerprint's collision probability to be at most $n^{-Ck-c_0}$, with sufficiently large constants $C,c_0$. Along the ideal execution for a fixed failure set, there are at most $O(nR\log n)$ tested recovery candidates. These candidates depend on the edge sampling and deterministic sparse-recovery sketches, not on the independent fingerprint randomness. A union bound over these tests and all $M_k\le n^k$ failure sets controls hallucinations separately from the progress estimate.

Each round uses $O(\log n)$ sampling levels and $O(k\log n)$ bits at each level. This proves
\[
 B_{\rm pub}=O\left(\min\{n,k\log^2n\log(en/k)\}\right).
\]

\subsection{Streaming initialization and randomness}
In a dynamic stream, the initial $2k$ neighbors need not be the first $2k$ in a fixed ordering. It suffices to recover any $2k$ genuine distinct neighbors per high-degree vertex. Maintain, at geometric sampling levels, an ordinary adjacency sparse-recovery sketch of capacity
\[
 d_0=C(k+\log n)
\]
and the exact sampled degree. If the sampled degree is at most $d_0$, recovery is certified by that degree; no fingerprint is necessary for this initialization. A vertex of degree at most $d_0$ is fully recovered at level zero. For a larger degree, a level with expected sample size between $d_0/8$ and $d_0/4$ has between $2k$ and $d_0$ sampled neighbors with probability $1-n^{-c-2}$, for sufficiently large $C$. A union bound handles all vertices. No uniformity of the returned $2k$ neighbors is needed for the large-component lemma.

The initialization costs $O((k+\log n)\log^2n)$ bits per vertex, which is absorbed because $k\log(en/k)\ge\log n$. Its randomness is independent of the subsequent merging-round randomness. Degree counts and linear recovery sketches support deletions in the stream.

The Bernoulli edge choices used in the main algorithm must also be reproducible from stored seeds. Use sufficiently high bounded independence to obtain the same exponentially small tail at mean $\Theta(k)$; $O(k)$-wise independence, with suitably chosen constants, suffices by the limited-independence concentration bounds of Schmidt--Siegel--Srinivasan~\cite{SSS95}. Polynomial hashing over a binary extension field supplies the geometric probabilities and uses $O(k\log n)$ seed bits per level and round. The initialization can likewise use $O(k+\log n)$-wise independence. The alternate fingerprint construction, corrected to use the whole field as described next, fits in the space budget. This establishes the claimed dynamic-streaming space, rather than accounting only for the sketches and omitting their random seeds.

\section{Specific issues in the supplied draft}
\paragraph{Off-by-one in the decision problem.}
Sections 5.1 and 5.2 enumerate $|F|\le k$ and require every remainder to be connected. Under the definition in Section 6, the correct condition is $|F|<k$. For example, the cycle $C_4$ is 2-connected, but deleting opposite vertices disconnects it. Keeping primitive support for $k$ failures is harmless; the decision decoder must use the correct threshold.

\paragraph{The alternate fingerprint proof uses a false distributional assertion.}
Section 5.3 samples entries of $\Phi$ uniformly from $\F_Q\setminus\{0\}$, then claims that independent linear projections are jointly uniform and independent. That is false on the punctured product distribution. For example, over an odd field, if $X,Y$ are independent uniform nonzero elements, $\Prb[X+Y=0]=1/(Q-1)$, not $1/Q$.

The simple repair is to sample every entry from the \emph{whole} field $\F_Q$. Conditional on $\Phi_F$, choose the kernel vectors as a deterministic function of $\Phi_F$. Their independence implies that the projection of each surviving column is uniform in $\F_Q^{c+1}$, with independent coordinates, and different surviving columns are independent. For a fixed nonzero boundary difference, each coordinate is a nonzero degree-two polynomial in its own independent set of projected variables. Schwartz--Zippel therefore gives total collision probability at most $(2/Q)^{c+1}$. Zero entries are entirely harmless for this fingerprint construction. This change is \emph{not} being imposed on the earlier edge-weight fingerprint construction, which has its own argument.

\paragraph{The merging proof is unfinished, and its quoted failure bound is vacuous for small $k$.}
The bound $2e^{-k/12}$ exceeds one for small $k$. Use a sufficiently large constant multiple of $k$ as recovery capacity and the progress proof in Section~\ref{sec:public}. Independence between blobs is neither available nor necessary. Do not condition on absence of future fingerprint collisions and then reuse unconditioned probabilities.

\paragraph{The private analysis needs a total-boundary event and a coupling.}
A component of $H-F$ has no boundary edges \emph{in $H-F$}. The intended boundary is in $G-F$. State the HKTZZ event as a bound on the \emph{total} number of edges joining different components of $H-F$, which automatically controls every component and supports the rank argument. Explain why surviving samples dominate a uniform sample in $G-F$, and respect the sample-size condition in the tail proof; Section~\ref{sec:private} provides these details.

\paragraph{A union of forests is not automatically a sparse certificate.}
The abstract promises a sparse graph, but merely unioning a forest for every failure set does not prove an $O(kn)$ or similar bound. Theorem~\ref{thm:sparsifier} supplies both the size guarantee and polynomial-time decoding.

\paragraph{The stronger abstract lower bound needs additional arguments.}
Theorem 4 is the promise lower bound $\Omega(\min\{\sqrt n,n/k\})$. It does not itself imply $\Omega(k+\sqrt n)$ for exact threshold decision. Theorem~\ref{thm:index} and Lemma~\ref{lem:join} supply the missing arguments. The proof of Theorem 4 itself appears sound in this audit, including the local-collision and Karpovsky--Milman steps.

\paragraph{Signs and parameter details.}
The displayed definition of $s(v)_t$ uses $\operatorname{sgn}(u-v)$, but Lemma 4.1 and its expansion use $\operatorname{sgn}(v-u)$. Already a single edge with $F=\varnothing,t=0$ exposes the mismatch. Fix one orientation throughout. Also require $Q>a_n^2$ explicitly, so distinct integer prime products remain distinct in the field; choosing a suitable $Q=\Theta(n^3)$ is compatible with this. The probability printed in Lemma 5.2 as $1-n^{-k+3}$ should not be used as a high-probability statement for small $k$.

\paragraph{Model and fingerprint quantifiers.}
The preliminaries define worst-case total communication, whereas the headline results use maximum per-vertex communication; define both. The fingerprint collision proof is pointwise for a fixed candidate independent of the fingerprint randomness. Arbitrarily adaptive purported vectors do not inherit that guarantee. In the applications, define ideal executions and union-bound their fixed candidates; for adaptive failure-set queries use the simultaneous all-$F$ guarantee.

\section{Optimal dimension of the generic linear deletion primitive}
The following is a generic linear-algebra fact, not a lower bound for all Boolean graph sketches.

\begin{proposition}
Let $d\ge1$. Suppose a linear sketch $A:\F^N\to\F^m$ supports the following task: from $Ax$ and any $F\subseteq[N]$ with $|F|\le k$, recover $x_{[N]\setminus F}$ whenever it has at most $d$ nonzero coordinates, with arbitrary values allowed on $F$. Then
\[
 m\ge\min\{N,2d+k\}.
\]
Over a sufficiently large field, a Vandermonde construction meets this dimension.
\end{proposition}
\begin{proof}
Any nonzero $z\in\ker A$ supported on at most $2d+k$ coordinates gives a contradiction. Choose $F$ to contain at most $k$ coordinates of $\supp z$, leaving between one and $2d$ nonzero coordinates outside $F$. Split those remaining coordinates into two sets of size at most $d$, and write $z=x-y$ so that both $x$ and $y$ are $d$-sparse outside $F$. The values on $F$ are unrestricted. Then $Ax=Ay$ but the required outputs differ. Thus every set of at most $2d+k$ columns must be linearly independent, which gives the dimension lower bound. With $2d+k$ Vandermonde rows, multiplying by the annihilating polynomial for $F$ leaves at least $2d$ moments, sufficient for sparse recovery. If $2d+k\ge N$, store $x$ itself.
\end{proof}

This explains why an additive deletion budget is natural, and permits removing the unnecessary final row in the draft's $2d+k+1$-row sparse-recovery matrix. It should not be advertised as an optimality theorem for the restricted $0/\pm1$ graph inputs without a separate argument.

\section{Checks, scope, and remaining research directions}
The accompanying \texttt{checks.py} is reproducible and uses no language-model agent. It checked:
\begin{itemize}[nosep]
\item all 24,057 combinations of a four-vertex graph, its sampled subgraph, up to two failed vertices, and three moment dimensions; this includes 22,593 sparse exact-rank cases and 88,842 moment-coordinate identities;
\item 264 small indexing gadgets, including even thresholds, padding, and universal-vertex lifting, checking both connectivity and independence of Alice's neighborhoods from Bob's index;
\item 592 graph/threshold cases for the polynomial sparsifier, checking the full component partition after every permitted failure set.
\end{itemize}
These computations supplement, and do not replace, the proofs.

The additional script \texttt{additional\_checks.py} checked 40,896 constant-field transformed matrices (including deliberately rank-deficient zero-matrix choices), 175,500 transformation identities, 36 nonzero cut polynomials on all 125 assignments over $\F_5$, the one-bit near-complete criterion on all 1,096 labeled graphs with three, four, or five vertices, and 5,324 edge-indispensability instances. It checked both rank soundness and the refinement property of the computed partition. Rank-deficient cases are expected and are not test failures.

The most substantive remaining targets are a polynomial-time \emph{one-sided} decoder for variable $k$, stronger private-randomness lower bounds, and removal of another logarithmic factor in the public upper bound. A related current result of Sawettamalya--Yu gives much smaller logarithmic overhead for \emph{edge}-connectivity certificates~\cite{SY26}; there is no justified black-box transfer of that theorem to vertex failures in this note. In particular, no $O(k\log^2n)$ bound with polynomial-time sparse-certificate decoding for general vertex-connectivity sketching is claimed here. The new rank decision protocol has a different tradeoff and must not be presented as such a certificate theorem.

\begin{thebibliography}{9}
\bibitem{SSZ09} G. B. Sorkin, A. Steger, and R. Zenklusen. \emph{A tight bound on the collection of edges in MSTs of induced subgraphs}. Journal of Combinatorial Theory, Series B 99(2), 428--435, 2009. arXiv:0705.2439. Theorem 6 bounds every $k$-constructible graph by $kn-\binom{k+1}{2}$ edges.
\bibitem{HKTZZ19} J. Holm, V. King, M. Thorup, O. Zamir, and U. Zwick. \emph{Random $k$-out subgraph leaves only $O(n/k)$ inter-component edges}. FOCS 2019. arXiv:1909.11147. In particular, Theorem 1.2, Corollary 2.22, and Section 3.
\bibitem{AS23} S. Assadi and V. Shah. \emph{Tight Bounds for Vertex Connectivity in Dynamic Streams}. SOSA 2023, pp. 213--227. arXiv:2211.04685. Appendix A reviews the greedy certificate construction and the Mader-based sparsity proof.
\bibitem{Yu21} H. Yu. \emph{Tight Distributed Sketching Lower Bound for Connectivity}. SODA 2021, pp. 1856--1873. arXiv:2007.12323.
\bibitem{RT26} P. Robinson and M. M. Tan. \emph{Polynomial Lower Bounds for Distributed Graph Sketching with Tiny Error: Connectivity and Spanning Tree Construction}. arXiv:2609.05795, 2026.
\bibitem{SY26} P. Sawettamalya and H. Yu. \emph{A (Very) Nearly Optimal Sketch for $k$-Edge Connectivity Certificates}. SOSA 2026, pp. 229--240. arXiv:2510.16336. Only the stated edge-connectivity result, not a vertex-connectivity adaptation, is invoked as a research direction here.
\bibitem{SSS95} J. P. Schmidt, A. Siegel, and A. Srinivasan. \emph{Chernoff--Hoeffding Bounds for Applications with Limited Independence}. SIAM Journal on Discrete Mathematics 8(2), 223--250, 1995.
\end{thebibliography}
\end{document}
